\section{Auswertung und Ergebnisse}
% Beispielinhalt: alles unterhalb der \section-Zeile kann gelöscht werden.

\blindtext

\begin{figure}[htbp]
  \centering
  \begin{subfigure}{0.45\textwidth}
    \includegraphics[width=\linewidth]{example-image-b}
    \caption{ohne Magnetfeld}
    \label{fig:spektrum-ohne}
  \end{subfigure}\hfill
  \begin{subfigure}{0.45\textwidth}
    \includegraphics[width=\linewidth]{example-image-c}
    \caption{mit Magnetfeld}
    \label{fig:spektrum-mit}
  \end{subfigure}
  \caption{Interferenzmuster der roten Cadmiumlinie. Lange Beschriftungen
    werden im Blocksatz gesetzt, kurze zentriert}
  \label{fig:spektren}
\end{figure}

\Cref{fig:spektren} zeigt die Interferenzmuster; in
\cref{fig:spektrum-mit} ist die Aufspaltung jeder Ordnung in drei Linien
zu erkennen. Aus den Abständen folgt mit \cref{eq:zeeman} die
Energieaufspaltung in \cref{tab:messwerte}.

\begin{table}[htbp]
  \caption{Theoretische und gemessene Energieaufspaltung $\Delta E$ in
    Abhängigkeit vom Magnetfeld $B$}
  \label{tab:messwerte}
  \centering
  \begin{tabular}{
      S[table-format=1.2]
      S[table-format=2.2]
      S[table-format=2.1(2)]}
    \toprule
    {$B$ / \unit{\tesla}} & {$\Delta E_\text{theo}$ / \unit{\micro\electronvolt}}
      & {$\Delta E_\text{exp}$ / \unit{\micro\electronvolt}} \\
    \midrule
    0.25 & 14.47 & 14.1(8)  \\
    0.50 & 28.94 & 29.5(8)  \\
    0.75 & 43.41 & 42.7(9)  \\
    1.00 & 57.88 & 58.6(10) \\
    \bottomrule
  \end{tabular}
\end{table}

\begin{figure}[htbp]
  \centering
  \includegraphics[width=0.7\textwidth]{example-image}
  \caption{Gemessene Energieaufspaltung über dem Magnetfeld mit linearer
    Anpassung}
  \label{fig:fit}
\end{figure}

Eine lineare Anpassung (\cref{fig:fit}) ergibt das Bohrsche Magneton zu
$\mu_\mathrm{B} = \qty{58.1(9)}{\micro\electronvolt\per\tesla}$.

\blindtext
